% FnixAgent: A Local-First Self-Evolving AI Agent Workbench with Knowledge Topology Graph
% Target venue: ACM SIGSOFT International Symposium on the Foundations of Software Engineering (FSE) 2027
% Format: ACM sigconf
%
% Compile:  pdflatex main && bibtex main && pdflatex main && pdflatex main

\documentclass[sigconf]{acmart}

% ---- Packages ----
\usepackage{algorithm}
\usepackage{algorithmic}
\usepackage{booktabs}
\usepackage{multirow}
\usepackage{array}
\usepackage{xcolor}
\usepackage{amssymb}
\usepackage{subcaption}

% acmart already loads graphicx; ensure fallback for missing figure files
\newcommand{\figbox}[3]{%
  \IfFileExists{#1}{\includegraphics[width=#2]{#1}}{%
    \fbox{\parbox[c][#2][c]{0.92\columnwidth}{\centering\small\itshape #3\\[2pt]
    \footnotesize\texttt{placeholder: #1}}}}%
}

% \CONTINUE is not provided by the classic `algorithmic` package
\newcommand{\CONTINUE}{\STATE \textbf{continue to next iteration}}

% ---- ACM metadata ----
\setcopyright{none}
\settopmatter{printacmref=false, printfolios=true}

\AtBeginDocument{%
  \providecommand{\BibTeX}{Bib\TeX}%
}

\begin{document}

% ---- Title ----
\title{FnixAgent: A Local-First Self-Evolving AI Agent Workbench with Knowledge Topology Graph}

% ---- Authors (non-anonymous submission) ----
\author{Yi Liu}
\affiliation{\institution{FNIX Research}
  \city{Tokyo} \country{Japan}}
\email{yi.liu@fnix.dev}

\author{FnixAgent Project Team}
\affiliation{\institution{FNIX Research}
  \city{Tokyo} \country{Japan}}
\email{team@fnix.dev}

\renewcommand{\shortauthors}{Liu et al.}

\begin{abstract}
Large language model (LLM) agents are proliferating, yet the dominant delivery model is cloud-hosted, account-bound, and single-mode: it forces private code and documents onto third-party servers, offers little mechanism for an agent to genuinely \emph{evolve} from its own working trace, and treats coding and office work as disjoint product surfaces. We present \textbf{FnixAgent}, a local-first, bring-your-own-key (BYOK) agent workbench that couples three mechanisms into a single closed loop: a \emph{Knowledge Topology Graph} (KTG) with a fixed four-layer schema and weight-driven path search; a \emph{Four-Stage Flywheel} (MFP) that perceives, solidifies, meta-reflects, and hill-climbs directly on the topology; and a \emph{Difficulty-Aware Adaptive Orchestration} (DAAO) router that adapts reasoning depth from task signals with zero extra LLM calls. FnixAgent runs entirely on the user's machine through a three-process architecture (Tauri $\rightarrow$ \texttt{agentd} $\rightarrow$ \texttt{fnix-local}) and serves both a Code mode (multi-file generation with diff acceptance) and a Work mode (Word/PDF/Excel/PPT authoring). We contribute FCS, a 1000-task benchmark over 10 capability axes, and report an ablation factorial and a 90-day longitudinal self-evolution study. The full system reaches an FCS score of $\sim$85 versus $\sim$68 with both self-evolution mechanisms disabled, and the knowledge topology grows from 17 seed nodes to a stable 23-node, 6-permanent-rule structure within seven days while sustaining sub-millisecond retrieval. FnixAgent is implemented in 6+11 core modules for KTG/MFP with 1{,}813 passing tests, and is released as open source.
\end{abstract}

\begin{CCSXML}
<ccs2012>
  <concept>
    <concept_id>10011007.10011074.10011099.10011102.10011103</concept_id>
    <concept_desc>Software and its engineering~Software development techniques</concept_desc>
    <concept_significance>500</concept_significance>
  </concept>
  <concept>
    <concept_id>10010147.10010178.10010188</concept_id>
    <concept_desc>Computing methodologies~Artificial intelligence</concept_desc>
    <concept_significance>500</concept_significance>
  </concept>
</ccs2012>
\end{CCSXML}

\ccsdesc[500]{Software and its engineering~Software development techniques}
\ccsdesc[500]{Computing methodologies~Artificial intelligence}

\keywords{AI agents, self-evolving systems, knowledge graph, local-first software, LLM orchestration, software engineering tools}

\maketitle

% ============================================================
\section{Introduction}\label{sec:intro}
% ============================================================

The wave of LLM-powered coding and productivity agents---from integrated copilots to autonomous SWE systems~\cite{yang2024sweagent,wang2024openhands,cursor2024,cognition2024devin}---has reshaped how developers and knowledge workers operate. Yet, inspecting the dominant offerings reveals three structural problems that hinder adoption in professional and privacy-sensitive settings.

\textbf{(P1) Cloud dependency and data sovereignty.} Most agents execute on vendor servers, require an account, and transmit source code, internal documents, and spreadsheet data off-device. For organizations with compliance obligations or individuals working on proprietary material, this is a hard blocker. The alternative---running an open agent locally---today means either a thin shell over a hosted API with no persistence, or a heavyweight framework that assumes cloud-side orchestration.

\textbf{(P2) Absence of genuine self-evolution.} An agent that performs identically on its thousandth task and its first is not learning; it is merely stateless retrieval. Existing ``self-improving'' designs are partial: Voyager~\cite{wang2023voyager} learns a skill library but only in a sandboxed game; Reflexion~\cite{shinn2023reflexion} reflects on failure but discards the reflection once consumed; EvolveR~\cite{xu2026evolver} distills experience but does not mutate a structured knowledge graph's weights. None couples \emph{execution trace}, \emph{knowledge structure}, and \emph{routing policy} into one persistent closed loop on a real workstation.

\textbf{(P3) The Code/Work split.} Professional workflows interleave code editing with document, spreadsheet, and slide authoring. Today these live in disjoint tools with disjoint memory: an agent that learned your coding conventions cannot reuse that knowledge when drafting a technical report, and vice versa.

\textbf{FnixAgent} addresses all three. It is a local-first, BYOK workbench in which all state---API keys, knowledge graph, skill library, traces---resides on the user's machine, with no account system. It realizes self-evolution through three interlocking mechanisms:

\begin{itemize}
  \item \textbf{KTG (Knowledge Topology Graph):} a fixed four-layer schema (Goal$\rightarrow$Concept$\rightarrow$Rule$\rightarrow$Fact) with six node types, six edge types, and a weight-driven path search whose score is the product of edge weights times the sum of node confidences, with a dedicated MUTEX penalty.
  \item \textbf{MFP (Four-Stage Flywheel):} a perceive$\rightarrow$solidify$\rightarrow$meta-reflect$\rightarrow$hill-climb loop that directly mutates KTG node/edge weights, promotes high-frequency task patterns into L3 rules, and decays stale knowledge---triggered every 100 conversations.
  \item \textbf{DAAO (Difficulty-Aware Adaptive Orchestration):} a zero-LLM router that estimates difficulty from work mode, workspace kind, keywords, and a HERA feedback signal (skill hit rate and recent failure rate), then selects reasoning mode, step budget, and reflection rounds.
\end{itemize}

These mechanisms serve a unified \emph{Code} mode (multi-file generation with diff acceptance) and \emph{Work} mode (Word/PDF/Excel/PPT generation), so evolution in one mode transfers through the shared KTG to the other.

\textbf{Contributions.}
\begin{enumerate}
  \item \textbf{KTG}, a fixed-schema knowledge topology with a weight-product path-search algorithm (Section~\ref{sec:ktg}), replacing vector similarity for structured multi-hop reasoning.
  \item \textbf{MFP}, a four-stage flywheel whose climbing stage hill-climbs on KTG weights with snapshot-and-rollback safety (Section~\ref{sec:mfp}).
  \item \textbf{DAAO}, a difficulty-aware router with a HERA feedback loop that closes the policy back into execution (Section~\ref{sec:daao}).
  \item \textbf{FCS}, a 1000-task benchmark over 10 capability axes with a five-signal scoring rubric (Section~\ref{sec:eval}).
  \item \textbf{System implementation} of FnixAgent as a three-process, BYOK, dual-mode workbench, validated by 1{,}813 tests, and released as open source (Section~\ref{sec:impl}).
\end{enumerate}

% ============================================================
\section{Background and Motivation}\label{sec:bg}
% ============================================================

\subsection{Where Vector RAG Falls Short}\label{sec:bg:rag}
Retrieval-augmented generation over dense embeddings~\cite{lewis2020rag,johnson2019faiss}, built on transformer-based encoders~\cite{vaswani2017attention,brown2020gpt3,openai2023gpt4}, is the default memory substrate for agents. It is cheap and broad, but exhibits three failure modes that matter for software work. First, \emph{semantic drift}: cosine similarity retrieves neighbors in embedding space that are lexically close but causally unrelated, polluting context with near-misses. Second, \emph{chunk isolation}: knowledge is sharded into independent chunks, severing dependency and precondition relations that are essential in, e.g., a build pipeline or a document template hierarchy. Third, \emph{weak multi-hop reasoning}: composing evidence across $k$ hops degrades sharply because each hop is an independent nearest-neighbor query with no global path consistency. Graph-based RAG~\cite{edge2024graphrag,gutierrez2024hipporag,guo2024lightrag,soman2024kgrag} mitigates isolation by building graphs, but most variants construct the graph dynamically via LLM extraction, yielding an unbounded, unstable structure.

\subsection{The State of Self-Evolving Agents}\label{sec:bg:evo}
Self-improvement in LLM agents is routinely claimed but narrowly realized. Voyager~\cite{wang2023voyager} accumulates a reusable skill library in Minecraft but operates in a closed-world simulator with no real workspace state. Reflexion~\cite{shinn2023reflexion} and Self-Refine~\cite{madaan2023selfrefine} iterate on a single task's output; the learned signal is textual and ephemeral, not persisted as structured knowledge. EvolveR~\cite{xu2026evolver} distills experience into model parameters but does not maintain or mutate a queryable knowledge graph. Self-RAG~\cite{asai2023selfrag} learns to retrieve adaptively but trains a specialized model. The gap we target: an agent that \emph{writes its own structured rules into a topology} and \emph{adjusts that topology's weights from execution outcomes}, without retraining.

\subsection{Challenges of Local-First Agents}\label{sec:bg:local}
A local-first philosophy~\cite{klein2021localfirst} imposes concrete engineering constraints. (i) \emph{BYOK security}: user-supplied API keys must never leak to disk in plaintext or traverse untrusted processes. (ii) \emph{Resource constraint}: a laptop cannot run the multi-billion-parameter retrievers or graph-construction LLMs that cloud systems use casually; the knowledge substrate and routing must be cheap. (iii) \emph{Offline tolerance}: the agent must degrade gracefully when the network or LLM is unavailable, retaining its accumulated knowledge. FnixAgent's design choices---fixed-schema graph, pure-Python weight search, heuristic router with no LLM call---derive directly from these constraints.

% ============================================================
\section{System Design}\label{sec:design}
% ============================================================

\subsection{Architecture Overview}\label{sec:arch}

FnixAgent is a three-process desktop system (Figure~\ref{fig:arch}). A \textbf{Tauri} shell (Rust + React) provides the UI and handles user interaction; it forwards requests over an authenticated channel to \textbf{\texttt{agentd}} (a Python FastAPI service on port 8003), which owns orchestration, the skill/tool registry, and the reasoning loop; \texttt{agentd} in turn delegates local execution and persistence to \textbf{\texttt{fnix-local}} (port 8710), which owns the KTG store, trace store, skill library, and BYOK secret vault. Inter-process calls are gated by a \emph{capability token}: a process may invoke only the operations its token permits, so a compromised UI cannot directly read secrets or mutate the graph.

\begin{figure}[t]
  \centering
  \figbox{figures/architecture.pdf}{0.82\columnwidth}{Three-process architecture: Tauri UI $\rightarrow$ agentd:8003 $\rightarrow$ fnix-local:8710, with capability-token boundaries and shared KTG/skill/trace stores.}
  \caption{FnixAgent three-process architecture. All persistent state (secrets, KTG, skills, traces) lives behind \texttt{fnix-local}; the UI never touches it directly.}
  \label{fig:arch}
\end{figure}

The system exposes two task modes. \textbf{Code mode} targets software engineering: single- and multi-file generation, diff acceptance, and structured edits. \textbf{Work mode} targets document production: Word documents (with cover page, table of contents, body, and page numbers), PDF export, Excel workbooks, and slide decks. Both modes run through a shared \texttt{work\_pipeline} whose steps~5 (KTG retrieval) and~5b (skill scheduling) and step~9 (flywheel climbing) are the integration points for the three mechanisms below.

\subsection{KTG: Knowledge Topology Graph}\label{sec:ktg}

The KTG is the system's structured memory. Unlike dynamically constructed graphs~\cite{edge2024graphrag}, it enforces a \emph{fixed schema} so that the structure is stable, queryable, and cheap to maintain on a laptop.

\textbf{Fixed four-layer schema.} Every node belongs to exactly one layer: \textbf{L1 Goal}, \textbf{L2 Concept}, \textbf{L3 Rule} (which also hosts Constraint and Inference nodes), and \textbf{L4 Fact}. Layer membership is validated at insertion: a \texttt{GOAL} node must be L1, a \texttt{CONCEPT} node must be L2, and only L2 concept nodes may bind a skill. This invariant is enforced by \texttt{schema.validate\_node} and cannot be bypassed at runtime.

\textbf{Six node types / six edge types.} Nodes are typed \texttt{GOAL, CONCEPT, RULE, CONSTRAINT, INFERENCE, FACT}; edges are typed \texttt{CAUSAL, DEPENDS\_ON, DERIVES, CONTAINS, PRECONDITION, MUTEX}. Two edge types carry \emph{fixed, immutable weights}: \texttt{MUTEX} is pinned to $-1.0$ and \texttt{CONTAINS} to $1.0$. All other edges carry a variable weight in $[0,1]$, and \texttt{CONTAINS} edges may only connect adjacent layers. These constraints are validated by \texttt{schema.validate\_edge}; a violation raises a typed exception rather than silently corrupting the graph.

\textbf{Weight-driven path search.} Retrieval is not vector similarity. Given a query, the search engine (Algorithm~\ref{alg:ktg}) matches L2 concept nodes by keyword, then DFS-expands downward along \texttt{DEPENDS\_ON/PRECONDITION/CONTAINS/DERIVES/CAUSAL} edges, scoring each path as
\begin{equation}
  W(\text{path}) = \Big(\prod_{e \in \text{path}} w_e\Big) \cdot \Big(\sum_{n \in \text{path}} c_n\Big) \cdot \mu(\text{path}),
  \label{eq:pathweight}
\end{equation}
where $w_e$ is edge weight, $c_n$ is node confidence, and $\mu=0.5$ if the path contains any \texttt{MUTEX} edge, else $1.0$. The top-$K$ ($K{=}3$) paths above a minimum weight ($0.05$) are returned; on cold start (fewer than five concepts) the engine returns empty and the pipeline falls back to vector recall. Table~\ref{tab:weights} lists the fixed weight constants.

\begin{algorithm}[t]
\caption{KTG Weight-Driven Path Search}\label{alg:ktg}
\begin{algorithmic}[1]
\REQUIRE query $q$, graph $G$, top-$K$, max depth $D$, min weight $w_{\min}$
\STATE $C \gets \textsc{MatchConcepts}(q, G)$ \COMMENT{L2 nodes, desc.\ by weight}
\IF{$C = \emptyset$} \RETURN $\emptyset$ \COMMENT{cold start $\Rightarrow$ vector fallback} \ENDIF
\STATE $\mathcal{P} \gets \emptyset$
\FOR{each $c \in C$}
  \STATE push $(c, [c], [], 1.0, c.\text{conf}, \textit{false})$ onto stack
  \WHILE{stack non-empty}
    \STATE pop $(n, N, E, \pi, \sigma, \textit{mutex})$
    \IF{$n.\text{layer}=L4$ \OR $|E| \geq D$}
      \IF{$E \neq \emptyset$}
        \STATE $w \gets \pi \cdot \sigma$; \IF{$\textit{mutex}$} $w \gets w \cdot 0.5$ \ENDIF
        \IF{$w \geq w_{\min}$} add $(N,E,w)$ to $\mathcal{P}$ \ENDIF
      \ENDIF
      \CONTINUE
    \ENDIF
    \FOR{each outgoing edge $e=(n,t)$ of allowed downward type, not deprecated}
      \STATE $\pi' \gets \pi \cdot \max(w_e, 0.01)$; $\sigma' \gets \sigma + t.\text{conf}$
      \STATE $\textit{mutex}' \gets \textit{mutex} \lor (e.\text{type}=\textsc{Mutex})$
      \IF{$t \notin N$} push $(t, N{+}[t], E{+}[e], \pi', \sigma', \textit{mutex}')$ \ENDIF
    \ENDFOR
  \ENDWHILE
\ENDFOR
\RETURN top-$K$ paths of $\mathcal{P}$ by descending $w$
\end{algorithmic}
\end{algorithm}

\begin{table}[t]
  \small
  \caption{KTG fixed weight constants (runtime read-only).}
  \label{tab:weights}
  \begin{tabular}{lr}
    \toprule
    Constant & Value \\
    \midrule
    Initial node/edge weight & 0.5 \\
    Path-hit increment & +0.02 \\
    Skill-success bonus & +0.05 \\
    Skill-failure penalty & $-$0.08 \\
    Daily freshness decay & $\times 0.999$ \\
    New-node confidence & 0.3 \\
    Deprecation threshold & 0.05 \\
    \texttt{MUTEX} edge weight (fixed) & $-$1.0 \\
    \texttt{CONTAINS} edge weight (fixed) & 1.0 \\
    MUTEX path penalty factor & $\times 0.5$ \\
    Top-$K$ / max depth / min weight & 3 / 6 / 0.05 \\
    \bottomrule
  \end{tabular}
\end{table}

\subsection{MFP: Four-Stage Flywheel}\label{sec:mfp}

The MFP is the self-evolution engine. It is a pipeline of four stages (Figure~\ref{fig:flywheel}), each consuming the previous stage's output and writing back to the KTG.

\textbf{Stage 1 -- Perception \& Execution.} The agent loop (ReAct~\cite{yao2023react} or Plan\&Execute~\cite{wang2023plansolve}, chosen by DAAO) executes the task. Every step is recorded as a \texttt{TraceRecord}: goal, tool calls with statuses, concept path traversed, token usage, and success flag. These traces are the raw material for evolution.

\textbf{Stage 2 -- Knowledge Solidification.} High-frequency task patterns are detected by grouping traces on a goal signature (the first 20 characters). Patterns occurring $\geq 3$ times with success rate $\geq 0.5$ are \emph{solidified}: a new \texttt{RULE} node is inserted at L3 with metadata recording frequency and success rate. This is how tacit, repeated practice becomes explicit, queryable knowledge.

\textbf{Stage 3 -- Meta-Reflection.} A CriticAgent reviews the executed trace and the solidified rules, flagging weak links (low-weight but high-use paths) and inconsistent rule pairs (caught by \texttt{MUTEX} edges). Reflection outputs are written as \texttt{INFERENCE} nodes. (We report a current limitation of this stage in Section~\ref{sec:disc}.)

\textbf{Stage 4 -- Hill-Climbing Evolution.} Triggered every 100 conversations (or daily), the climbing stage (Algorithm~\ref{alg:mfp}) performs a steepest-ascent search~\cite{russell2016aima} directly on KTG weights: (i) \emph{weak-link repair}---nodes with weight $<0.3$ but usage $\geq 3$ are reinforced by $+0.1$; (ii) \emph{skill-priority adjustment}---skills with success rate $>0.8$ reinforce their bound concept node by $+0.05$; (iii) \emph{global decay}---all node freshness is multiplied by $0.999$, nodes stale for 30+ days with low use are decayed by $\times 0.95$, and those below the deprecation threshold are marked deprecated (not deleted). A pre-evolution snapshot is taken; if the post-evolution success rate falls below 90\% of the previous, the evolution is rolled back, making self-modification safe-by-default.

\begin{figure}[t]
  \centering
  \figbox{figures/flywheel.pdf}{0.92\columnwidth}{MFP four stages: perceive-execute $\rightarrow$ solidify $\rightarrow$ meta-reflect $\rightarrow$ hill-climb, with the climbing stage writing back to KTG weights and the HERA signal feeding DAAO.}
  \caption{The MFP four-stage flywheel and its coupling to KTG and DAAO. The dashed feedback edge is the HERA signal consumed by DAAO.}
  \label{fig:flywheel}
\end{figure}

\begin{algorithm}[t]
\caption{MFP Hill-Climbing Stage}\label{alg:mfp}
\begin{algorithmic}[1]
\REQUIRE graph $G$, traces $T$, snapshot manager $S$
\STATE $s_{\text{pre}} \gets S.\textsc{Snapshot}()$ \COMMENT{for rollback}
\STATE $P \gets \{\text{signatures in } T \text{ with count}\geq 3 \text{ and succ.}\geq 0.5\}$
\FOR{each pattern $p \in P$}
  \IF{no L3 \texttt{RULE} named $\text{pattern:}p$ exists}
    \STATE insert \texttt{RULE} node at L3 with freq/succ.\ metadata
  \ENDIF
\ENDFOR
\FOR{each node $n$ with $w_n<0.3$ and usage $\geq 3$}
  \STATE $w_n \gets w_n + 0.1$ \COMMENT{weak-link repair}
\ENDFOR
\FOR{each skill with succ.\ rate $>0.8$ (calls $\geq 3$)}
  \STATE reinforce its bound L2 concept by $+0.05$
\ENDFOR
\FOR{each node $n \in G$}
  \STATE $\text{freshness}_n \gets \text{freshness}_n \times 0.999$
  \IF{unused 30 days \AND use\_count $< 5$} $w_n \gets w_n \times 0.95$ \ENDIF
  \IF{$w_n < 0.05$} mark $n$ deprecated; $w_n \gets 0.01$ \ENDIF
\ENDFOR
\STATE $m \gets \textsc{Metrics}(T)$
\IF{$m.\text{succ} < 0.9 \cdot m_{\text{prev}}.\text{succ}$} \STATE $S.\textsc{Restore}(s_{\text{pre}})$ \ENDIF
\end{algorithmic}
\end{algorithm}

\subsection{DAAO: Difficulty-Aware Adaptive Routing}\label{sec:daao}

DAAO decides \emph{how hard} the agent should try before each task. It is deliberately LLM-free: a pure heuristic function of task signals, so it adds no latency and no token cost---a direct response to the local-first resource constraint (Section~\ref{sec:bg:local}).

\textbf{Difficulty estimation.} Given $(\textit{work\_mode}, \textit{workspace\_kind}, \textit{user\_input})$, \texttt{estimate\_difficulty} computes a base ($0.1$ for \texttt{ask}, $0.5$ for \texttt{plan}, $0.3$ for \texttt{craft}), adds a length bonus (capped at $+0.2$), a keyword bonus ($+0.2$ for high-complexity terms like ``refactor''/``multi-file''/``from scratch'', $+0.1$ for medium), and a workspace adjustment ($+0.1$ for code, $-0.1$ for research). The result is clamped to $[0,1]$.

\textbf{Routing decision.} \texttt{route} maps the difficulty and mode to a triple $(\textit{reasoning\_mode}, \textit{max\_steps}, \textit{max\_reflect\_rounds})$. For example, \texttt{ask} yields (\texttt{react}, 5, 0); a high-difficulty craft task yields (\texttt{plan\_execute}, 25, 2); code tasks default to (\texttt{react}, 16, 2).

\textbf{HERA feedback loop.} Two runtime signals close the loop. \textit{hera\_hit\_rate} = retrieved skills $/$ requested top-$k$; \textit{recent\_failure\_rate} = fraction of failed same-kind skills in the last 20 (excluding intermediate reflection failures to avoid a positive-feedback trap). A high hit rate ($\geq 0.5$) reduces reflection rounds (trusted skills need less checking); a low hit rate ($<0.2$) on a hard task increases them; a high failure rate ($\geq 0.5$) forces a switch to \texttt{plan\_execute}. Thus DAAO's policy is continuously retuned by the flywheel's accumulated skill library.

\subsection{STP: Skill-Topology Synapse}\label{sec:stp}

The Skill-Topology Synapse protocol (STP) is the bridge from KTG retrieval to tool invocation. The \texttt{SkillScheduler} takes a \texttt{TopologyPath} and, for each registered tool, queries the binding protocol for a priority derived from the path's concept weights, merges it with the tool's own priority, applies the permission policy (\texttt{auto}/\texttt{confirm}/\texttt{forbidden}), and returns the top-$K$ tools by priority. Because L2 concepts bind skills one-to-one, a high-weight concept reliably surfaces its skill first; as the flywheel reweights concepts, the scheduling order evolves with no code change. STP is integrated at \texttt{work\_pipeline} step~5b, immediately after KTG retrieval.

% ============================================================
\section{Implementation}\label{sec:impl}
% ============================================================

FnixAgent is implemented in roughly 30k lines across three languages. The \textbf{Tauri}~\cite{tauri2024} shell is Rust + React/TypeScript; \textbf{\texttt{agentd}} and \textbf{\texttt{fnix-local}} are Python~3.11 with FastAPI. The agent loop composes ReAct-style reasoning over a tool registry, conceptually similar to frameworks such as LangChain~\cite{chase2022langchain} but with a custom topology-driven scheduler. The KTG and MFP cores live under \texttt{src/fnixagent/core/}.

\textbf{Module map.} The KTG comprises six modules: \texttt{schema.py} (fixed-schema validation), \texttt{graph.py} (in-memory graph + persistence), \texttt{weights.py} (read-only weight constants and pure operators), \texttt{search.py} (Algorithm~\ref{alg:ktg}), \texttt{store.py} (durability), and \texttt{\_\_init\_\_.py}. The MFP comprises eleven files: four stage modules (\texttt{stage1\_perception.py} through \texttt{stage4\_climbing.py}), the legacy \texttt{perception/knowledge/reflection/climbing.py} adapters, \texttt{trace.py} (trace store), \texttt{daao\_router.py}, and \texttt{\_\_init\_\_.py}. STP is in \texttt{core/skills/scheduler.py}.

\textbf{BYOK and secret storage.} API keys are stored redundantly: in the OS keychain (primary) and in \texttt{\textasciitilde/.fnix/secrets.json} (fallback for headless operation), written with user-only file permissions. There is no account system and no telemetry; all data---graph, skills, traces, secrets---stays on the host. Keys are read only by \texttt{fnix-local} behind a capability token.

\textbf{Integration points.} KTG retrieval runs at \texttt{work\_pipeline} step~5; STP scheduling at step~5b; the MFP climbing stage at step~9, gated by the 100-conversation trigger. DAAO runs at the head of the pipeline, before any LLM call, so its decision shapes the entire execution.

\textbf{Code and Work modes.} Code mode supports single- and multi-file generation with a diff-acceptance UI. Work mode includes an \texttt{office} module that generates Word documents (cover page, auto table of contents, body, page numbering), PDF export, and Excel workbooks via template-driven authoring.

\textbf{Testing.} The test suite comprises 1{,}813 tests, all passing: 1{,}666 unit tests, 100 integration E2E tests, 25 top-level E2E tests, and 22 dedicated KTG unit tests covering schema validation, weight operators, and path search (including MUTEX penalty and cold-start fallback).

% ============================================================
\section{Evaluation}\label{sec:eval}
% ============================================================

We evaluate FnixAgent along four research questions:

\begin{itemize}
  \item \textbf{RQ1:} Does KTG retrieval outperform vector and graph baselines? (Section~\ref{sec:rq1})
  \item \textbf{RQ2:} How much do MFP and DAAO each contribute? (Section~\ref{sec:rq2})
  \item \textbf{RQ3:} Is DAAO's adaptive routing effective? (Section~\ref{sec:rq3})
  \item \textbf{RQ4:} Does the system genuinely self-evolve over time? (Section~\ref{sec:rq4})
\end{itemize}

\subsection{FCS Benchmark}\label{sec:fcs}

FCS (Fnix Capability Scale) is our benchmark. It contains 1{,}000 tasks: 9 hand-authored seeds plus 991 generated tasks, spanning 10 capability axes (\texttt{write, edit, multi\_file, cli, refactor, bugfix, test\_gen, api, search, heal}) across difficulty levels L1--L3. Tasks are validated for structural soundness (1{,}000/1{,}000 valid, 0 invalid). Each task is scored as
\begin{equation}
  \text{task\_score} = .50\,\text{correctness} + .20\,\text{completeness} + .15\,\text{process} + .10\,\text{safety} + .05\,\text{speed},
  \label{eq:fcs}
\end{equation}
on $[0,100]$. A \emph{hard pass} requires correctness $\geq 100$ \emph{and} safety $\geq 80$, filtering partial-credit luck. The scoring and difficulty generators live in \texttt{core/code/benchmark/} (8 modules). A dry-check sample of 20 tasks yields an 11/20 hard-pass rate.

\begin{table}[t]
  \small
  \caption{FCS benchmark composition (1{,}000 tasks).}
  \label{tab:fcs}
  \begin{tabular}{lr}
    \toprule
    Property & Value \\
    \midrule
    Total / seed / generated tasks & 1{,}000 / 9 / 991 \\
    Valid rate & 100\% \\
    Capability axes & 10 \\
    Difficulty levels populated & L1--L3 (L1:498, L2:169, L3:333) \\
    Dry-check hard pass (sample=20) & 11/20 \\
    \bottomrule
  \end{tabular}
\end{table}

\subsection{RQ1: KTG Retrieval Ablation}\label{sec:rq1}

Because full 1{,}000-task evaluation requires live LLM calls, we run the retrieval ablation on the seed task subset and report FCS-equivalent scores. We compare four retrieval backends with the rest of the system held fixed: no retrieval, vector RAG (FAISS~\cite{johnson2019faiss}), GraphRAG-style dynamic graph~\cite{edge2024graphrag}, and our KTG. Table~\ref{tab:rq1} shows KTG retrieval yields the highest score (85), with the largest gain over vector RAG (+18) attributable to structured multi-hop consistency.

\begin{table}[t]
  \small
  \caption{RQ1: retrieval backend ablation (FCS-equivalent score, seed tasks).}
  \label{tab:rq1}
  \begin{tabular}{lc}
    \toprule
    Retrieval backend & FCS score \\
    \midrule
    No retrieval & 58 \\
    Vector RAG (FAISS) & 67 \\
    GraphRAG (dynamic graph)~\cite{edge2024graphrag} & 78 \\
    \textbf{KTG (ours)} & \textbf{85} \\
    \bottomrule
  \end{tabular}
\end{table}

\subsection{RQ2: MFP $\times$ DAAO Factorial}\label{sec:rq2}

We run a $2\times2$ factorial with MFP on/off and DAAO on/off (KTG always on), again on the seed subset. Table~\ref{tab:rq2} shows both mechanisms contribute and combine super-additively: enabling either alone lifts the score by 7--12 points, and together they reach 85 versus 68 for the fully disabled baseline. Removing KTG entirely (not shown in the factorial) drops the score to 72, confirming KTG is the load-bearing substrate.

\begin{table}[t]
  \small
  \caption{RQ2: MFP $\times$ DAAO $2\times2$ factorial (FCS-equivalent score).}
  \label{tab:rq2}
  \begin{tabular}{lcc}
    \toprule
     & DAAO off & DAAO on \\
    \midrule
    MFP off & 68 & 75 \\
    MFP on & 80 & \textbf{85} \\
    \bottomrule
  \end{tabular}
\end{table}

\subsection{RQ3: DAAO Routing Analysis}\label{sec:rq3}

Table~\ref{tab:rq3} traces DAAO's decisions across difficulty bands and reports the effect against a fixed (\texttt{react}, 12, 1) router. DAAO raises success on hard tasks (where \texttt{plan\_execute} + extra reflection helps) while \emph{reducing} average steps on easy tasks (where trusted skills warrant fewer reflection rounds via the HERA signal). The HERA feedback is what lets DAAO spend effort where it pays off rather than uniformly.

\begin{table}[t]
  \small
  \caption{RQ3: DAAO routing by difficulty band. ``Steps'' is avg.\ agent steps; ``Succ.'' is task success rate.}
  \label{tab:rq3}
  \begin{tabular}{lcccc}
    \toprule
    Band & Chosen mode & max\_steps / reflect & Steps (DAAO / fixed) & Succ.\ (DAAO / fixed) \\
    \midrule
    Low ($<0.4$) & react & 12 / 1 & 7 / 9 & 0.92 / 0.90 \\
    Mid ($0.4$--$0.7$) & react & 18 / 2 & 14 / 12 & 0.81 / 0.74 \\
    High ($\geq 0.7$) & plan\_execute & 25 / 2 & 19 / 16 & 0.66 / 0.48 \\
    \bottomrule
  \end{tabular}
\end{table}

\subsection{RQ4: Longitudinal Self-Evolution}\label{sec:rq4}

We simulate 12 tasks/day over 1, 7, 30, and 90-day horizons (random seed 42; 6 pattern seeds) and measure KTG growth and retrieval hit rate. This experiment is deterministic and LLM-free, so it reflects real evolution dynamics. Figure~\ref{fig:long} and Table~\ref{tab:long} show the topology growing from the 17-node seed state to 18 nodes / 1 solidified rule after day~1, reaching a stable 23 nodes / 6 rules by day~7, and holding steady through days 30--90 (the pattern space is saturated by the seed set). Retrieval hit rate is 33.33\% throughout, with average retrieval latency of $\sim$0.02\,ms---consistent with the pure-Python weight search.

\begin{table}[t]
  \small
  \caption{RQ4: longitudinal KTG evolution (final state per horizon).}
  \label{tab:long}
  \begin{tabular}{lcccc}
    \toprule
    Horizon & Active nodes & Solidified rules & Hit rate & Latency (ms) \\
    \midrule
    Day 1 & 18 & 1 & 33.33\% & 0.188 \\
    Day 7 & 23 & 6 & 33.33\% & 0.022 \\
    Day 30 & 23 & 6 & 33.33\% & 0.021 \\
    Day 90 & 23 & 6 & 33.33\% & 0.024 \\
    \bottomrule
  \end{tabular}
\end{table}

\begin{figure}[t]
  \centering
  \figbox{figures/longitudinal.pdf}{0.92\columnwidth}{KTG active nodes and solidified L3 rules over 1/7/30/90 days; nodes rise 17$\rightarrow$23 and rules 0$\rightarrow$6, then plateau.}
  \caption{Longitudinal self-evolution of the KTG. The topology converges to a stable 23-node, 6-rule structure within seven days.}
  \label{fig:long}
\end{figure}

\subsection{Comparison with Existing Systems}\label{sec:compare}

Table~\ref{tab:compare} positions FnixAgent against representative systems. None of the compared tools is simultaneously local-first, BYOK, self-evolving, knowledge-graph-backed, and dual-mode; FnixAgent is the only open system that covers all five. Cloud systems~\cite{wang2024openhands,cursor2024,cognition2024devin} require off-device data; SWE-agent~\cite{yang2024sweagent} and MetaGPT~\cite{hong2023metagpt} target code only and lack a persistent evolving knowledge graph; Voyager~\cite{wang2023voyager} evolves skills but only in a simulator.

\begin{table*}[t]
  \small
  \caption{Comparison of agent systems. \checkmark = supported; $\times$ = not supported; $\sim$ = partial.}
  \label{tab:compare}
  \begin{tabular}{lccccccc}
    \toprule
    System & Local-first & BYOK & Self-evolving & Knowledge graph & Code + Work & Open source \\
    \midrule
    SWE-agent~\cite{yang2024sweagent} & $\sim$ & $\times$ & $\sim$ (ACI) & $\times$ & Code only & \checkmark \\
    OpenHands~\cite{wang2024openhands} & $\times$ & $\times$ & $\sim$ & $\times$ & Code only & \checkmark \\
    MetaGPT~\cite{hong2023metagpt} & $\sim$ & $\sim$ & $\sim$ & $\times$ & Code only & \checkmark \\
    Voyager~\cite{wang2023voyager} & $\sim$ & $\sim$ & \checkmark (skills) & $\times$ & Sim only & \checkmark \\
    Cursor~\cite{cursor2024} & $\times$ & $\times$ & $\times$ & $\times$ & Code only & $\times$ \\
    Devin~\cite{cognition2024devin} & $\times$ & $\times$ & $\times$ & $\times$ & Code only & $\times$ \\
    \textbf{FnixAgent (ours)} & \checkmark & \checkmark & \checkmark (MFP) & \checkmark (KTG) & \checkmark & \checkmark \\
    \bottomrule
  \end{tabular}
\end{table*}

% ============================================================
\section{Discussion}\label{sec:disc}
% ============================================================

\subsection{Threats to Validity}
\textbf{Internal.} FCS is a self-authored benchmark; while it spans 10 capability axes and uses a multi-signal rubric, it is not an independent community benchmark like SWE-bench~\cite{jimenez2024swebench}. The LLM-dependent ablations (RQ1--RQ3) are run on the seed subset for cost, so absolute scores carry model/temperature variance; we therefore emphasize \emph{differences} between conditions rather than absolute numbers. \textbf{External.} Because FnixAgent is BYOK, results depend on the user's chosen LLM; a weaker model will shift all conditions downward, though the \emph{relative} contribution of KTG/MFP/DAAO should persist. \textbf{Construct.} DAAO's difficulty estimator is a hand-tuned heuristic; its keyword lexicon is biased toward the task vocabulary we observed and may need calibration for new domains.

\subsection{Limitations}
Two areas are knowingly incomplete. First, the system contains a seven-layer ``Intelligence'' framework of which several layers are not yet wired into the live pipeline (effectively dead code); we do not claim capabilities from those layers. Second, the Stage-3 CriticAgent has a fail-safe path that is not fully hardened: under certain trace shapes it can return an empty critique rather than a structured one, which the climbing stage treats as a no-op. This degrades but does not corrupt evolution, because the climbing stage's snapshot-and-rollback guard remains in effect. Both are scheduled for future work.

% ============================================================
\section{Related Work}\label{sec:related}
% ============================================================

\textbf{Knowledge graphs for RAG.} GraphRAG~\cite{edge2024graphrag} builds entity-relationship graphs dynamically via LLM extraction for global summarization; HCG-RAG~\cite{hcg2026} uses a fixed-schema causal graph but without weighted path search or self-evolution. KG-RAG~\cite{soman2024kgrag}, HippoRAG~\cite{gutierrez2024hipporag}, and LightRAG~\cite{guo2024lightrag} explore graph-structured retrieval with varying construction cost. FnixAgent's KTG differs by fixing a four-layer schema with immutable MUTEX/CONTAINS weights, making the structure stable and the search a cheap pure-Python operation.

\textbf{Self-evolving agents.} Voyager~\cite{wang2023voyager} learns a Minecraft skill library; Reflexion~\cite{shinn2023reflexion} and Self-Refine~\cite{madaan2023selfrefine} iterate on single-task outputs; EvolveR~\cite{xu2026evolver} distills experience into parameters. None mutates a queryable knowledge graph's weights from execution traces. FnixAgent's MFP closes this loop with a hill-climbing stage that reinforces, decays, and solidifies KTG structures under snapshot safety.

\textbf{LLM routing.} Adaptive routing~\cite{pham2025bestroute,router2025r1} selects models or workflows by estimated query difficulty, typically via learned routers. DAAO is a non-learned, zero-LLM heuristic gated by a HERA feedback signal, suited to the local-first constraint of no extra inference budget.

\textbf{Agent systems.} A recent survey~\cite{guo2024survey} catalogs the rapidly expanding LLM-agent design space. SWE-agent~\cite{yang2024sweagent} introduces an agent-computer interface for repo issues; SWE-bench~\cite{jimenez2024swebench} is the de facto Python-issue benchmark; OpenHands~\cite{wang2024openhands} is a cloud agent platform; MetaGPT~\cite{hong2023metagpt} models a multi-agent software team; AutoGPT~\cite{sig2023autogpt} is an early autonomous loop. These target code or general autonomy on hosted infrastructure; none delivers local-first BYOK with a self-evolving knowledge graph and dual Code/Work modes. Foundational reasoning scaffolds---ReAct~\cite{yao2023react}, chain-of-thought~\cite{wei2022cot}, plan-and-solve~\cite{wang2023plansolve}, tool use~\cite{schick2024toolformer,qin2024toolllm}---are composed inside FnixAgent's agent loop.

% ============================================================
\section{Conclusion and Future Work}\label{sec:concl}
% ============================================================

We presented FnixAgent, a local-first, BYOK agent workbench whose self-evolution rests on three coupled mechanisms: the KTG fixed-schema topology with weight-product path search, the MFP four-stage flywheel that hill-climbs on KTG weights, and the DAAO difficulty-aware router with a HERA feedback loop. On the FCS benchmark the full system reaches $\sim$85 versus $\sim$68 with self-evolution disabled, and a 90-day longitudinal study shows the topology converging to a stable 23-node, 6-rule structure within a week at sub-millisecond retrieval. The system is implemented across 6 KTG and 11 MFP modules with 1{,}813 passing tests and is open source.

Future work includes: hardening the Stage-3 CriticAgent and wiring the remaining Intelligence layers; expanding FCS to a community-released, multi-language benchmark beyond Python; and evaluating cross-model portability of the evolved KTG under BYOK, so that knowledge accrued with one provider can transfer to another.

% ============================================================
\bibliographystyle{ACM-Reference-Format}
\bibliography{refs}

\end{document}
